\documentclass[11pt]{article}

\usepackage[T1]{fontenc}
\usepackage{newtxtext}
\usepackage{newtxmath}
\usepackage{amsmath}
\usepackage{amssymb}
\usepackage[margin=1in]{geometry}
\usepackage{natbib}
\usepackage{enumitem}
\usepackage{booktabs}
\usepackage{graphicx}
\usepackage{float}
\usepackage{xcolor}
\usepackage{hyperref}
\hypersetup{hidelinks}

\setlength{\parindent}{0pt}
\setlength{\parskip}{0.5em}
\setlist[itemize]{leftmargin=1.3em,itemsep=0.2em,topsep=0.2em}
\setlist[enumerate]{leftmargin=1.5em,itemsep=0.2em,topsep=0.2em}

\newcommand{\tbd}[1]{\textbf{[TBD: #1]}}

\begin{document}

{\Large\bfseries Historical Transportability of Private-Capital SDFs\par}
{\normalsize Which historical cash flows are relevant for pricing current funds?\par}
\today

\begin{abstract}
Private-capital SDF estimation faces a three-way data constraint. There are few
vintage years, older vintages may transport poorly to current funds, and the
most recent vintages have not yet produced the realized cash flows needed for
appraisal-free estimation. We propose a relevance-weighted hybrid estimator
that treats the historical information set itself as a model-selection object.
Seasoned vintages contribute realized-cash-flow NPV moments; unseasoned
vintages can contribute short-horizon NAV Euler errors; and historical
vintages are weighted by temporal proximity and similarity in credit spreads,
interest rates, and public-market valuations. Expanding, rolling, decaying,
and state-conditioned SDFs are compared on a common NAV-free criterion. An
ex-post transportability matrix measures how well an SDF estimated in one
historical window prices earlier and later realized vintages. A forward
pseudo-real-time exercise then jointly selects the SDF specification, the
historical weighting scheme, and whether unseasoned Euler moments improve
subsequent realized-cash-flow pricing relative to excluding young funds. The
paper does not claim that thirty vintage years can cleanly distinguish a
stable law from every form of temporal instability. It asks the empirically
sharper question: which use of the available history travels best out of
sample?
\end{abstract}

\section{Introduction}

How much historical private-capital data should be used to estimate an SDF
for a fund raised today? Under a stable and sufficiently mixing data-generating
process, older vintages add information and a full-history estimator becomes
increasingly precise. Under temporal instability, however, the same estimator
may converge precisely to a historical mixture that is no longer the relevant
pricing law for current funds. More funds within a vintage can reveal that
vintage's pricing moment more accurately, but they cannot create additional
independent economic histories.

This creates a three-way empirical constraint. First, the time-series dimension
is intrinsically thin: even a long private-capital database contains only a few
dozen vintage years, and adjacent vintages live through overlapping calendar
histories. Second, increasing the sample by reaching further into the past
requires a transport assumption: that pricing relations estimated on older
funds remain relevant for funds raised under different financing, valuation,
and exit conditions. Third, restricting estimation to fully realized cash
flows removes precisely the recent vintages that are most informative about
the current environment. Conventional resampling can reduce fund noise inside
the observed sample, but it cannot create economic time or resolve this
maturity delay.

The private-equity pricing literature has developed several ways to recover
risk and performance from irregular, non-traded cash flows.
\citet{DriessenLinPhalippou2012} estimate factor exposures and abnormal
performance directly from fund cash flows, while
\citet{AngChenGoetzmannPhalippou2018} recover a time series of private-equity
returns from limited-partner cash flows. \citet{KortewegNagel2016} use a
generalized public-market-equivalent approach to risk-adjust venture-capital
cash flows. More recent work refines cash-flow benchmarking and SDF-based
performance measurement \citep{GredilGriffithsStucke2023,KortewegNagel2025}.
These approaches differ in their maintained pricing restrictions and treatment
of illiquidity, but they largely take the historical estimation sample as
given. Our paper instead asks which part of that sample is informative for a
current pricing problem: how well does an estimated private-capital SDF
transport across economic time?

We call the empirical ability of an SDF estimated in one historical window to
price cash flows from another window \emph{historical transportability}. A
stable pricing relation should travel across windows; a temporally unstable
one should exhibit loss that changes with the distance and economic difference
between estimation and evaluation periods. The concept does not require a
literal regime classification. Smooth drift, isolated breaks, changes in the
support of economic states, and changes in fund composition can all appear as
failures of transportability. In an environment this short, a global test of
``stability versus instability'' is unlikely to have useful power against the
many plausible alternatives. We therefore follow the forecasting-under-
instability logic of evaluating local predictive performance rather than
making a structural-break test the paper's central result
\citep{PesaranTimmermann2007,Rossi2021}.

Our proposed solution is a \emph{forward-selected relevance-weighted hybrid
SDF}. It has three components. First, candidate SDFs are estimated under a
small, pre-specified sequence of historical weighting schemes: the full
expanding sample, hard rolling windows, exponential time decay, and a
time-and-state kernel. Second, recent economic similarity is measured with
three observable state variables fixed before evaluating fund performance:
credit spreads, the interest-rate level, and public-market valuations. Third,
recent unseasoned vintages enter through short-horizon NAV Euler restrictions,
but their inclusion is not maintained by assumption: the hybrid estimator is
compared with an otherwise identical estimator that excludes them. Every
choice is judged on later realized cash flows.

The empirical design has two parts. An ex-post transportability matrix uses
fully realized vintages to show how pricing laws travel forward and backward
across historical windows. Its backward entries are structural diagnostics,
not feasible forecasts. A pseudo-real-time exercise then recreates historical
information dates, estimates each candidate using the information available
at that origin, and scores it on cash flows realized later by the then-young
funds. This second exercise provides the directional evidence required to
choose a current estimator.

The paper makes four contributions.
\begin{enumerate}
\item It formulates the private-capital data problem as a joint constraint of
      few economic histories, potentially decaying historical relevance, and
      delayed cash-flow outcomes.
\item It introduces historical transportability matrices and signed-distance
      loss curves that show where and how far a private-capital SDF travels.
\item It proposes a relevance-weighted SDF that combines simple temporal
      windows with low-dimensional economic-state similarity, while enforcing
      an effective-sample-size and identification constraint.
\item It evaluates the informational value of unseasoned NAV Euler moments on
      a common, later-realized cash-flow target rather than on NAV fit.
\end{enumerate}

\section{Methodology}

\subsection{The Private-Capital Data Trilemma}

The empirical problem has three connected parts:
\begin{enumerate}
\item \textbf{Few vintages.} The relevant time-series sample contains only a
      few dozen vintage years, and overlapping fund lives make its effective
      temporal dimension smaller than its vintage count.
\item \textbf{Historical relevance.} Adding older vintages reduces sampling
      variance only if their pricing relations transport to the target period.
      Under temporal instability, a precisely estimated full-history SDF can
      be precisely wrong for a current fund.
\item \textbf{Delayed outcomes.} The newest vintages contain the most current
      economic information but do not yet reveal their full realized cash
      flows. Discarding them mechanically makes a current estimator historical.
\end{enumerate}

These are not three independent nuisances. Shortening the estimation window
raises temporal relevance but increases sampling uncertainty and concentrates
the sample in unseasoned funds. Requiring realized cash flows pushes the
effective window backward. Reaching further into history restores moment count
at the cost of a stronger transport assumption. The empirical objective is
therefore not to eliminate any one component in isolation, but to find the
combination that produces the lowest forward realized-cash-flow loss.

\subsubsection{What the data can and cannot establish}

Thirty vintage years do not support a clean omnibus decision between a
time-invariant SDF and every form of temporal instability. A failure to reject
stability may reflect low power, while a local episode of poor performance need
not identify a permanent break. The paper therefore reports economically
interpretable loss differences, transport curves, and model sets across a
small pre-specified design space. Statistical tests are companions to this
evidence, not the sole basis for declaring a regime change. In the terminology
of the forecasting literature, local forecast performance is the target;
parameter instability is one possible explanation for its variation
\citep{Rossi2021}.

\subsubsection{The cost of ignoring recent vintages}

Discarding unseasoned vintages is clean only in a narrow measurement sense. It
also changes the estimand by restricting the current SDF to funds raised far
enough in the past to have substantially realized. Because buyout persistence
and the composition of the industry have changed over time, this restriction
cannot automatically be treated as innocuous \citep{HarrisJenkinsonKaplanStucke2023}.
The exclusion estimator is therefore an essential benchmark, not the default
solution. Whether short-horizon Euler information improves forward pricing is
an empirical question.

\subsection{Pricing Moments}

\subsubsection{Seasoned vintages}

For seasoned fund $i$ of vintage $v$, let the realized pricing error per dollar
committed be
\[
g_i^S(\theta;M)
=
\frac{1}{K_i}\sum_{t=0}^{T_i}\Psi_{0,t}(\theta;M)CF_{t,i},
\]
and let $\epsilon_v^S(\theta;M)$ be its pooled vintage mean. The time-zero
capital call enters undiscounted. The realized-cash-flow block is the common
estimation and evaluation currency throughout the paper.

\subsubsection{Unseasoned vintages}

For an unseasoned fund, the full NPV moment is unavailable without valuing the
residual claim. The baseline bridge uses the short-horizon NAV Euler error
\[
e^H_{t,i}(\theta;M)
=
M_{t-1,t}(\theta;M)R^H_{t,i}-1,
\]
where $R^H_{t,i}$ is the quarterly Horizon-IRR return accounting for interim
calls and distributions. Fund-quarter errors are aggregated and mapped into a
vintage moment $\epsilon_v^U(\theta;M)$ using a pre-specified duration
normalization. They are never discounted back to inception and never enter the
comparison score. Reported NAVs are noisy appraisals, so this block is treated
as a candidate information bridge rather than a traded-return restriction that
must hold exactly. State-space corrected NAVs are a robustness extension, not
the baseline \citep{brown2023nowcasting}.

The main unseasoned-treatment comparison is deliberately parsimonious:
\begin{enumerate}
\item \emph{Exclude}: estimate only with seasoned realized-cash-flow moments;
\item \emph{Hybrid}: add short-horizon NAV Euler moments for unseasoned
      vintages.
\end{enumerate}
Both are evaluated on the identical subsequently realized cash flows. A more
complicated measurement model earns a place in the main design only if the
simple comparison demonstrates economically meaningful value from recent NAV
information.

\subsection{The Relevance-Weighted Hybrid SDF}

At target date $T$, let $z_v$ collect the economic state associated with
vintage $v$. The baseline state vector contains three pre-performance
variables: a credit spread, the interest-rate level, and a public-market
valuation ratio. The variables are measured at or immediately before vintage
formation and standardized using the estimation sample only.

For historical vintage $v$, define the unnormalized relevance weight
\[
w_{v,T}(h,b)
=
\exp\!\left[-\log(2)\frac{T-v}{h}\right]
\exp\!\left[-\frac{1}{2b^2}
  (z_v-z_T)'(z_v-z_T)\right],
\]
where $h$ is a time half-life and $b$ is a common state-similarity bandwidth.
The simple common bandwidth avoids estimating three separate state-sensitivity
parameters from a few dozen vintages. Setting $h=\infty$ removes time decay;
setting $b=\infty$ removes state conditioning. Hard rolling windows are
obtained by replacing the time-decay component with an indicator.

Let $q_v$ be a pre-specified reliability weight and normalize $w_{v,T}q_v$ to
sum to one within the relevant estimation block. For candidate unseasoned
treatment $A\in\{\mathrm{Exclude},\mathrm{Hybrid}\}$, estimate
\[
\widehat\theta_T(M,h,b,A)
=
\arg\min_{\theta\in\Theta_M}
\left\{
\sum_{v\in\mathcal V_T^S}\widetilde w_{v,T}
  \epsilon_v^S(\theta;M)^2
+
\mathbf 1\{A=\mathrm{Hybrid}\}
\sum_{v\in\mathcal V_T^U}\widetilde w_{v,T}
  \epsilon_v^U(\theta;M)^2
\right\}.
\]

The design sequence moves from transparent to reasonably sophisticated while
remaining feasible in a thin vintage panel:
\begin{enumerate}
\item \textbf{Expanding}: equal weight on the full available history;
\item \textbf{Rolling}: equal weight inside a fixed recent-vintage window;
\item \textbf{Decay}: exponential time weights indexed by half-life $h$;
\item \textbf{Time--state}: exponential decay combined with similarity in
      credit spreads, interest rates, and public valuations.
\end{enumerate}

A regularized time-varying parameter path of the form
\[
\min_{\{\theta_v\}}
\sum_v\epsilon_v(\theta_v)^2
+\lambda\sum_{v>1}\|\theta_v-\theta_{v-1}\|_2^2
\]
is the principal sophistication check \citep{CuiFengHong2024}. It is not the
headline estimator unless it produces a clear forward improvement over the
lower-dimensional weighting schemes.

Every candidate must satisfy two feasibility restrictions before it can be
scored: the weighted moment Jacobian must be sufficiently well conditioned,
and the weight concentration must leave an effective vintage count
\[
V_{\mathrm{eff}}(T)=\frac{1}{\sum_v\widetilde w_{v,T}^2}
\]
comfortably above the number of SDF parameters. This prevents an apparently
local estimator from being driven by one historically similar vintage.

\subsection{Historical Transportability}

\subsubsection{The transportability matrix}

Let $s$ index a historical estimation window and $t$ an excluded, fully
realized evaluation vintage. Define the NAV-free transport loss
\[
L_{s\rightarrow t}(M,h,b)
=
\widehat U_t^{\mathrm{real}}\!\left(
  \widehat\theta_s(M,h,b,\mathrm{Exclude})
\right),
\]
where the evaluation score uses only realized fund cash flows. The exclusion
version isolates historical transportability from the separate question of
how young-fund NAV information enters. Stacking these losses over estimation
windows and evaluation vintages yields a
\emph{historical transportability matrix}. Its rows are window SDFs, its
columns are evaluation vintages, and each entry measures how well a pricing
law travels across economic time.

Aggregating by signed temporal distance $d=t-s$ gives the transport loss curve
\[
\mathcal T_{M,h}(d)
=
\frac{1}{N_d}\sum_s L_{s\rightarrow s+d}(M,h,b).
\]
Positive distances measure forward transportability; negative distances
measure backward transportability. Only the former has a real-time forecasting
interpretation. Backward transportability is a diagnostic of structural
stability, not a feasible investment rule.

The matrix is descriptive by design. A flat surface is consistent with broad
historical transportability; loss increasing in $|d|$ is consistent with
temporal localization; and isolated high-loss columns indicate historically
unusual target periods. None of these patterns alone proves a particular
structural-break model.

\subsubsection{Calendar-time overlap}

Adjacent vintage moments share public-factor histories because their funds
remain active for many years. We therefore measure pairwise calendar-time
overlap rather than treating vintage labels as independent by construction.
Let $a_{v,t}$ denote vintage $v$'s discounted-cash-flow or Jacobian exposure to
calendar period $t$. Define
\[
O_{vv'}
=
\frac{\sum_t a_{v,t}a_{v',t}}
{\sqrt{\sum_ta_{v,t}^2}\sqrt{\sum_ta_{v',t}^2}}.
\]
The overlap matrix is reported alongside the transportability matrix and is
used to purge near-duplicate evaluation windows where feasible. Because
complete separation would eliminate most of the data, overlap is quantified
and stress-tested rather than treated as a binary admissibility standard.

\subsection{Unseasoned Vintages as Delayed Outcomes}

At historical information date $\tau$, let $A$ denote a bridge for unseasoned
vintages and estimate
\[
\widehat\theta_{\tau,h,A}(M)
=
\arg\min_{\theta\in\Theta_M}Q_{\tau,h,A}(\theta;M)
\]
using only fund information dated no later than $\tau$. Candidate bridges are
compared on the identical realized-cash-flow target after the then-unseasoned
vintages have seasoned:
\[
L_{\tau\rightarrow\tau+\ell}(M,h,A)
=
\frac{1}{|\mathcal V^U_\tau\cap\mathcal V^S_{\tau+\ell}|}
\sum_{v\in\mathcal V^U_\tau\cap\mathcal V^S_{\tau+\ell}}
\widehat U_v^{\mathrm{real}}
\!\left(\widehat\theta_{\tau,h,A}(M)\right).
\]

The empirical comparison therefore selects a triplet: the SDF specification,
the historical weighting rule, and the unseasoned-vintage treatment. Exclusion
remains in the race as the benchmark measuring the cost of delayed outcomes;
the short-horizon Euler block is the proposed simple bridge. The paper presents
the configuration with the lowest forward loss as its preferred empirical
solution, subject to the feasibility and stability diagnostics above.

\subsection{Evaluation Design}

\textbf{Ex-post transportability.} Estimate local SDFs on fully observed
historical windows and score them on earlier and later seasoned vintages. This
exercise is fully NAV-free and isolates temporal transportability, but it does
not claim a real-time interpretation when later information is used to estimate
an SDF that is scored on an earlier period.

\textbf{Pseudo-real-time delayed-outcome validation.} Reconstruct a sequence of
historical information dates, estimate each window/bridge combination using
the information fields available at the corresponding date, and score it on
subsequently realized fund cash flows. A literal real-time reconstruction would
also require historical database snapshots, contemporaneous coverage records,
and knowledge of which observations were later backfilled. Imposing that ideal
would leave essentially no usable evaluation sample. We therefore adopt the
pragmatic pseudo-real-time standard feasible in the data and state its residual
look-ahead limitation explicitly.\footnote{The fully reconstructed real-time
ideal is conceptually useful even when empirically infeasible: it identifies
the remaining source of look-ahead without pretending that an empty sample is
a superior research design.}

At each origin, tuning uses only earlier admissible origins. The outer origin
is reserved for evaluation. With few origins, the tuning grid is intentionally
coarse and common across SDF specifications. The empirical result is reported
as a loss surface and a set of near-optimal configurations, alongside the
single preferred configuration, so that small differences are not given a
false winner-takes-all interpretation.

\subsection{Forward Selection Criterion}

Let $\mathcal O$ denote the admissible pseudo-real-time origins and
$\mathcal H_\tau$ the vintages that were unseasoned at $\tau$ and later become
eligible for realized-cash-flow scoring. For design
$c=(M,h,b,A)$, define
\[
\widehat R(c)
=
\frac{1}{|\mathcal O|}
\sum_{\tau\in\mathcal O}
\frac{1}{|\mathcal H_\tau|}
\sum_{v\in\mathcal H_\tau}
\widehat U_v^{\mathrm{real}}
\!\left(\widehat\theta_\tau(c)\right).
\]
The preferred design minimizes $\widehat R(c)$. All designs face the same
origins and realized evaluation funds, so comparisons are paired. The paper
emphasizes effect sizes, origin-by-origin stability, and sensitivity to the
coarse tuning grid. With only a few dozen vintage years, formal rejections of
equal predictive loss are interpreted cautiously.

\subsection{Inference and Reporting}

Model comparisons use paired loss differences because every admissible design
is evaluated on the same origin--vintage cells. Losses are first averaged
within origin so that origins with many funds or newly seasoned vintages do not
dominate mechanically. We report the mean paired difference, its distribution
across origins, and block-bootstrap uncertainty that preserves adjacent-origin
dependence. Parameter uncertainty for a selected SDF is assessed with vintage-
level resampling and calendar-overlap blocks. Given the short time series and
the additional uncertainty created by model selection, confidence intervals
and predictive-loss tests are interpreted as measures of precision rather than
as a binary stability verdict.

\section{Data and Empirical Design}

This section maps the empirical inputs directly into the moments, state
weights, and forward loss defined above. The ordering is deliberate: we first
construct an appraisal-free realized-cash-flow target, then define what was
observable at each historical origin, and only then compare alternative uses of
the available history.

\subsection{Private-Capital Cash-Flow Data}

The fund data are drawn from \tbd{data provider and extraction date}. The unit
of observation is a dated limited-partner cash flow for a \tbd{buyout / broader
private-capital} fund. Required fields are fund identifier, vintage year,
commitment, dated capital calls, dated distributions, reported NAVs, strategy,
geography, and --- where available --- general-partner identifier. Cash flows
and NAVs are measured net of management fees and carried interest from the
limited partner's perspective and converted to \tbd{currency and real/nominal
units}. We exclude funds with internally inconsistent cash-flow histories,
missing commitments, nonpositive commitments, or implausible scale ratios
under a pre-specified cleaning rule. The analysis retains all admissible funds
within a vintage; it does not create pseudo-vintages by resampling funds.

The first data table should make the thin time-series dimension visible rather
than reporting only a large fund count:

\begin{table}[htbp]
\centering
\caption{Private-capital sample by vintage period}
\label{tab:sample}
\small
\begin{tabular}{lrrrrr}
\toprule
Vintage period & Vintages & Funds & Commitments & Funds/vintage & Realized share \\
\midrule
\tbd{early period}  & & & & & \\
\tbd{middle period} & & & & & \\
\tbd{recent period} & & & & & \\
\midrule
Full sample & & & & & \\
\bottomrule
\end{tabular}
\begin{minipage}{0.94\textwidth}
\footnotesize \emph{Notes:} Commitments are reported in \tbd{units}. Funds per
vintage is the median across vintage years. Realized share is the fraction of
funds satisfying the baseline seasoning rule at the final sample date.
\end{minipage}
\end{table}

\subsection{Seasoning and the Realized Evaluation Sample}

At information date $\tau$, a fund is seasoned when its residual value share
\[
s_{i,\tau}
=
\frac{\mathrm{RVPI}_{i,\tau}}
{\mathrm{DPI}_{i,\tau}+\mathrm{RVPI}_{i,\tau}}
\]
is at most $\delta$. The proposed baseline is $\delta=0.15$, with $0.10$ and
$0.20$ as pre-specified robustness thresholds. A fund that fails this rule may
contribute to the hybrid short-horizon Euler block but never to the realized
evaluation loss at that date. The realized target includes only calls and
distributions observed by the evaluation date; the residual NAV is used to
determine eligibility but is neither priced nor added as a terminal cash flow.
Thus, $\delta$ trades bounded terminal truncation for sample size without
fitting the score to NAV; $\delta=0$ is the strict robustness check.

We report, for every pseudo-real-time origin, the number of available vintages,
seasoned and unseasoned funds, the fraction of commitment represented by the
realized evaluation set, and the later vintages that can eventually be scored.
This origin-by-origin coverage table is essential: it reveals whether an
apparent forecasting gain comes from a genuinely broader information set or
from a small and unusual set of funds that happens to mature early.

\subsection{Public Factors and Candidate SDFs}

Public-market returns and the risk-free rate are obtained from \tbd{source} at
the same frequency as the cash-flow discounting exercise. The main comparison
uses a deliberately small SDF set: a risk-free benchmark, a one-factor public-
equity SDF, a levered public-equity specification, and at most one established
multifactor alternative. This keeps the empirical question focused on the
historical information set rather than a factor zoo. Factor definitions,
currency, return convention, and the mapping from quarterly factors to dated
cash flows are fixed before inspecting forward pricing losses.

\subsection{Economic-State Variables}

The state vector $z_v$ contains one variable for each of three economic
dimensions: a corporate credit spread, the interest-rate level, and a public-
equity valuation ratio. The proposed baseline definitions are the high-yield
option-adjusted spread, a long-term nominal government yield, and the cyclically
adjusted price--earnings ratio \tbd{or a price--dividend ratio}. To avoid using
fund outcomes to define similarity, each variable is measured as its average
over the 12 months ending immediately before the vintage year. At each origin,
levels are standardized using only the available estimation vintages. We also
report the historical support of $z_T$; a target outside that support is an
extrapolation problem and should not receive a falsely precise local estimate.

\subsection{Pseudo-Real-Time Origins}

Annual origins begin at the first date for which all candidate SDFs satisfy the
minimum effective-vintage and Jacobian conditions. At each origin, cash flows,
NAVs, factors, and state variables are truncated to their admissible information
date. Hyperparameters are chosen using only earlier origins, and the current
origin is used once for evaluation. The final extract does not reconstruct
historical database coverage or later backfills; this pragmatic limitation is
disclosed and tested by excluding observations with suspiciously late first
reporting where such timestamps exist. The stronger literal real-time standard
would otherwise reduce the usable evaluation sample to essentially zero.

\subsection{Pre-Specified Design Grid}

The grid is intentionally coarse. Its final endpoints are fixed after a
descriptive feasibility audit, but before comparing out-of-sample losses.

\begin{table}[htbp]
\centering
\caption{Pre-specified empirical design}
\label{tab:design}
\small
\begin{tabular}{p{0.23\textwidth}p{0.29\textwidth}p{0.39\textwidth}}
\toprule
Object & Baseline candidates & Restriction or robustness \\
\midrule
SDF & Risk-free; public equity; levered equity; one multifactor model
    & Common factor set at every origin \\
Historical weighting & Expanding; rolling; exponential decay; time--state
    & Same candidate sequence for every SDF \\
Rolling length & 8, 12, 16 vintages
    & Retained only when identification is adequate \\
Time half-life & 4, 8, 12 vintage years
    & $h=\infty$ reproduces no decay \\
State bandwidth & 0.75, 1.5, 3 standardized units
    & $b=\infty$ reproduces time-only weighting \\
Young-fund treatment & Exclude; short-horizon NAV Euler hybrid
    & Common later realized-cash-flow target \\
Seasoning threshold & Residual value share at most 15\%
    & 10\%, 20\%, and zero residual value \\
Feasibility & $V_{\mathrm{eff}}\geq\max\{8,3\dim(\theta)\}$
    & Jacobian conditioning reported separately \\
\bottomrule
\end{tabular}
\end{table}

\subsection{Empirical Workflow}

The empirical analysis proceeds in the same order as the methodology:
\begin{enumerate}
\item construct seasoned realized-cash-flow moments and document calendar-time
      overlap across vintages;
\item estimate fully realized local SDFs and form the ex-post historical
      transportability matrix;
\item reconstruct pseudo-real-time origins and calculate feasible expanding,
      rolling, decaying, and time--state estimates;
\item select tuning parameters using earlier origins only;
\item score all candidates on the same later realized-cash-flow outcomes and
      compare exclusion with the short-horizon NAV Euler hybrid; and
\item re-estimate the selected specification at the final information date and
      report both its current SDF parameters and the historical weights that
      support them.
\end{enumerate}

\section{Empirical Results}

\subsection{The Effective Historical Sample}

Table~\ref{tab:sample} first documents the nominal and effective sample size.
A vintage-by-calendar-year cash-flow heat map shows the long overlap of fund
lives, while the overlap matrix $O$ summarizes common factor exposure. We
report the eigenvalue-based effective temporal rank and effective vintage
counts under each candidate weighting rule. This section establishes how much
independent economic history is actually available before asking whether old
history remains relevant.

\subsection{Ex-Post Historical Transportability}

The first main figure is the historical transportability matrix for the
baseline SDF. Rows identify estimation windows, columns identify excluded
realized vintages, and colors represent realized pricing loss. Signed-distance
loss curves then summarize forward and backward transport. The empirical
questions are whether loss rises with temporal distance, whether the pattern is
asymmetric, and whether particular evaluation vintages are difficult for every
historical window.

\subsection{How Much History Travels Forward?}

The main model-comparison table reports paired forward loss for the expanding,
rolling, decay, and time--state estimators on common origin--vintage cells. For
each design we report average loss, the difference relative to expanding
history, origin-level uncertainty, $V_{\mathrm{eff}}$, and the fraction of
origins passing identification checks. This table supplies the paper's primary
answer: whether using less, or selectively weighted, history improves pricing
of later realized private-capital cash flows.

\subsection{Does Economic-State Similarity Add Information?}

We next hold the time-weighting rule fixed and compare time-only with time--state
weights. Weight plots show which distant vintages are recovered because their
credit spreads, interest rates, or public valuations resemble the target.
Leave-one-state-variable-out results reveal whether improvement is broad-based
or driven by a single state proxy. A state-conditioned estimator is preferred
only if it improves forward loss without collapsing the effective vintage
count.

\subsection{Do Unseasoned Vintages Improve the Current SDF?}

For every admissible origin, we compare the exclusion and short-horizon NAV
Euler estimators using identical later realized cash flows. The main table
reports paired loss differences overall and by the age of the young funds at
the estimation origin. We also relate the hybrid's gain to its NAV share and
subsequent NAV revisions. This distinguishes useful early information from a
mechanical fit to smoothed reported values.

\subsection{The Preferred SDF and Current Risk-Adjusted Performance}

The final results section re-estimates the forward-selected design at the
current information date. We report SDF parameters, implied factor exposures,
risk-adjusted performance, the weights assigned to every historical vintage,
and sensitivity to all configurations whose forward loss is statistically and
economically close to the minimum. The conventional expanding-history estimate
is shown alongside it so that the effect of historical transportability is
transparent.

\section{Robustness and Method Validation}

\subsection{Alternative Sample and Moment Choices}

Robustness exercises vary the seasoning threshold, exclude residual NAV
entirely, remove the financial-crisis and pandemic vintages one at a time,
alter fund-size reliability weights, and restrict the sample by strategy or
geography where coverage permits. Results are also reported with vintage-
equal weighting so that a few data-rich vintage years cannot determine the
conclusion.

\subsection{Alternative Windows, States, and SDFs}

We vary the coarse window grid, use a real rather than nominal long rate,
replace the valuation proxy, winsorize state distances, and compare the common
bandwidth with a strongly regularized variable-specific bandwidth. The smooth
time-varying parameter path is retained as a sophistication check. Alternative
SDFs matter only insofar as the transportability result survives reasonable
pricing specifications.

\subsection{Placebos}

Two placebos discipline the interpretation of state similarity. First, state
vectors are permuted across vintages while their time ordering is retained.
Second, the target state is replaced by a future state that was unavailable at
the origin. Genuine transport gains should disappear under the first placebo;
the second quantifies how much an inadvertently forward-looking state
construction could exaggerate performance.

\subsection{Compact Method Validation}

A standalone simulation section is not required for the paper's main empirical
claim. The estimator is selected by observed forward realized-cash-flow loss,
and a large menu of simulated breaks would not establish which process
generated the single historical private-capital sample. A compact appendix
calibration can nevertheless provide a useful diagnostic. It should match the
empirical vintage count, cash-flow-age profile, and calendar overlap under only
two designs --- a stable SDF and a smoothly drifting SDF --- each with and
without appraisal smoothing. Its limited purpose is to verify that the
selection procedure does not mechanically prefer short windows, state
conditioning, or the NAV bridge when the pricing law is stable. The empirical
transport and delayed-outcome exercises remain the evidence of interest.

\clearpage
\section{Conclusion}

Private-capital SDF estimation cannot solve a thin and delayed economic history
by resampling funds more aggressively. The relevant choice is which historical
vintages to transport to the target period and whether partial information from
recent funds improves that transport. The proposed estimator makes both choices
empirical: it weights history by time and observable economic similarity,
admits recent NAV information only through a short-horizon bridge, and selects
the complete design by its ability to price later realized cash flows. The
result is a current SDF supported by the part of history that has demonstrated
the greatest forward relevance, rather than by the longest history available.

\bibliographystyle{plainnat}
\bibliography{resampling_refs}

\end{document}
